% =====================================================
% commands.tex — Variables propias de cada enunciado
% Editar estos campos al inicio para cada ejercicio
% =====================================================

% --- Identificación del ejercicio ---
\newcommand{\tituloEjercicio}{Ejercicio: [TÍTULO]}
\newcommand{\subtituloEjercicio}{[Subtítulo descriptivo]}
\newcommand{\codigoEjercicio}{TEKLA-XXX}
\newcommand{\nivelEjercicio}{Avanzado}
\newcommand{\duracionEjercicio}{[X] horas}

% --- Datos del curso ---
\newcommand{\nombreCurso}{[Nombre del curso]}
\newcommand{\autorEnunciado}{[Docente]}
\newcommand{\fechaEjercicio}{\today}

% --- Versión del software ---
\newcommand{\versionTekla}{Tekla Structures 2022}
\newcommand{\entornoTekla}{[Environment / Configuration]}

% --- Atajos para terminología recurrente ---
\newcommand{\CC}{\tk{Component Catalog}}
\newcommand{\AC}{\tk{Applications \& Components}}
\newcommand{\OB}{\tk{Object Browser}}
\newcommand{\OR}{\tk{Object Representation}}
\newcommand{\OG}{\tk{Object Group}}
\newcommand{\SF}{\tk{Selection Filter}}
\newcommand{\VF}{\tk{View Filter}}
